% =====================================================================
\section{Шпаргалка по C++ для Arduino}
% =====================================================================
Здесь собраны конструкции языка, которые встречаются во всех уроках. Это не учебник C++, а справка:
открывайте её, когда забыли, как записывается условие или цикл.

\subsection*{Типы данных}

\begin{center}\small
\begin{tabularx}{\textwidth}{@{}l l X@{}}
\toprule
Тип & Диапазон на ESP32-S3 & Когда использовать \\
\midrule
\code{int} & от $-2\,147\,483\,648$ до $2\,147\,483\,647$ & обычные целые числа, счётчики \\
\code{uint8\_t} & 0\ldots255 & байт, яркость, значение цвета \\
\code{uint32\_t} & 0\ldots4\,294\,967\,295 & время в \code{millis()}, большие неотрицательные числа \\
\code{float} & около 7 значащих цифр & дробные числа (температура, напряжение) \\
\code{bool} & \code{true} или \code{false} & флаги, состояния «включено/выключено» \\
\code{char} & один символ \code{'A'} & отдельные символы \\
\code{String} & строка произвольной длины & текст в Arduino (удобно, но использует память) \\
\bottomrule
\end{tabularx}
\end{center}

\begin{Verbatim}[fontsize=\small,bgcolor=codebg]
int count = 0;                  // целое
float temp = 23.5;              // дробное
bool ledOn = false;             // логическое
uint32_t lastTime = 0;          // для millis()
String name = "ESP32";          // строка
\end{Verbatim}

\subsection*{Константы}

Значение, которое не должно меняться, объявляют с \code{const}. Компилятор не даст случайно изменить
его, а имена принято писать ЗАГЛАВНЫМИ буквами.

\begin{Verbatim}[fontsize=\small,bgcolor=codebg]
const int LED_PIN = 4;          // номер вывода
constexpr float VREF = 3.3;     // известно ещё на этапе компиляции
#define BUTTON_PIN 5            // старый способ: просто подстановка текста
\end{Verbatim}

\subsection*{Операторы}

\begin{center}\small
\begin{tabularx}{\textwidth}{@{}l X@{}}
\toprule
Операторы & Назначение \\
\midrule
\code{+ - * / \%} & арифметика; \code{\%} --- остаток от деления \\
\code{== != < > <= >=} & сравнение (результат \code{true} или \code{false}) \\
\code{\&\& || !} & «и», «или», «не» \\
\code{= += -= ++ --} & присваивание, увеличение и уменьшение \\
\code{\& | \^{} << >>} & побитовые операции (маски, сдвиги) \\
\bottomrule
\end{tabularx}
\end{center}

\begin{Verbatim}[fontsize=\small,bgcolor=codebg]
int a = 7 % 3;                  // 1 (остаток)
bool ok = (a > 0) && (a < 5);   // true
count++;                        // то же, что count = count + 1
mask |= (1 << 3);               // включить бит номер 3
\end{Verbatim}

\subsection*{Условия: if / else и switch}

\begin{Verbatim}[fontsize=\small,bgcolor=codebg]
if (temp > 30) {
  digitalWrite(FAN_PIN, HIGH);
} else if (temp < 25) {
  digitalWrite(FAN_PIN, LOW);
} else {
  // между 25 и 30 ничего не меняем
}

switch (mode) {                 // выбор по значению
  case 0: Serial.println("выкл"); break;
  case 1: Serial.println("авто"); break;
  default: Serial.println("?");
}
\end{Verbatim}

\subsection*{Циклы: for и while}

\begin{Verbatim}[fontsize=\small,bgcolor=codebg]
for (int i = 0; i < 5; i++) {   // пять повторов: i = 0..4
  Serial.println(i);
}

uint32_t start = millis();
while (!ready && millis() - start < 5000) {  // ждём не дольше 5 с
  delay(10);
}
\end{Verbatim}

Цикл \code{while (true)} без выхода и без тайм-аута в учебных программах не пишите: плата «повиснет» (урок 14).

\subsection*{Функции}

Функция даёт имя повторяющемуся действию. У неё есть тип результата, имя и список параметров.
Тип \code{void} означает «ничего не возвращает».

\begin{Verbatim}[fontsize=\small,bgcolor=codebg]
int clampPercent(int x) {       // вернуть число, ограниченное 0..100
  if (x < 0) return 0;
  if (x > 100) return 100;
  return x;
}

void blinkOnce(int pin, int ms) {   // ничего не возвращает
  digitalWrite(pin, HIGH);
  delay(ms);
  digitalWrite(pin, LOW);
}

int level = clampPercent(130);  // level = 100
\end{Verbatim}

\subsection*{Массивы}

\begin{Verbatim}[fontsize=\small,bgcolor=codebg]
int pins[3] = {4, 5, 6};        // три элемента: индексы 0, 1, 2
int n = sizeof(pins) / sizeof(pins[0]);  // число элементов = 3

for (int i = 0; i < n; i++) {
  pinMode(pins[i], OUTPUT);
}
// pins[3] -- ошибка: выход за границы массива!
\end{Verbatim}

\subsection*{Структуры}

Структура объединяет несколько значений под одним именем: удобно для пакета данных (урок 12) или
для описания датчика.

\begin{Verbatim}[fontsize=\small,bgcolor=codebg]
struct Reading {
  float temp;
  float humidity;
  uint32_t time;
};

Reading r = {22.5, 40.0, 0};
r.temp = 23.1;                  // обращение к полю через точку
Serial.println(r.temp);
\end{Verbatim}

\subsection*{Область видимости и static}

Переменная, объявленная внутри функции, существует только в ней и при каждом вызове создаётся заново.
Переменная вне функций (глобальная) видна всем. Модификатор \code{static} внутри функции сохраняет значение
между вызовами, но оставляет имя видимым только в этой функции.

\begin{Verbatim}[fontsize=\small,bgcolor=codebg]
int total = 0;                  // глобальная: видна везде

void step() {
  int local = 0;                // создаётся заново при каждом вызове
  static int calls = 0;         // помнит значение между вызовами
  local++;
  calls++;
  total += 10;
}
\end{Verbatim}

\subsection*{Типичные ошибки}

\begin{center}\small
\begin{tabularx}{\textwidth}{@{}>{\raggedright\arraybackslash}p{4.6cm} X@{}}
\toprule
Ошибка & Как правильно \\
\midrule
\code{if (x = 5)} --- присваивание вместо сравнения & \code{if (x == 5)}: сравнение записывают двумя знаками \\
Деление целых: \code{int a = 1 / 2;} даёт 0 & хотя бы одно число дробное: \code{1.0 / 2} или \code{(float)1 / 2} \\
Переполнение: \code{uint8\_t b = 250; b += 10;} даёт 4 & выбирайте тип с запасом (\code{int}, \code{uint32\_t}) \\
Сравнение времени: \code{millis() > t + 1000} ломается через 49 суток & пишите вычитание: \code{millis() - t > 1000} \\
Выход за границы массива & индексы от 0 до размер$-1$; проверяйте индекс \\
Точка с запятой сразу после условия: \code{if (x > 0);} & блок после условия ставят без неё: \code{if (x > 0) \{ \ldots \}} \\
Сравнение строк через \code{==} с \code{char*} & для \code{String} можно \code{==}, для \code{char*} --- \code{strcmp()} \\
\bottomrule
\end{tabularx}
\end{center}

\begin{Verbatim}[fontsize=\small,bgcolor=codebg]
int half = 7 / 2;               // 3, а не 3.5: целочисленное деление
float f = 7 / 2.0;              // 3.5
uint8_t b = 250;
b += 10;                        // переполнение: теперь b = 4
if (millis() - lastTime >= 1000) { lastTime = millis(); }  // верно
\end{Verbatim}
